% =====================================================================
%  BioTerm-Bench v1 -- ECCB conference poster, A0 portrait (841 x 1189 mm).
%
%  Build:   make poster        (runs xelatex twice; tectonic also works)
%
%  Figures are poster-scale redraws produced by
%  scripts/make_poster_figures.py from the same CSV/JSON receipts that back
%  the manuscript, and both tables are \input from paper/generated/, so the
%  poster cannot drift from main.tex: rerun scripts/refresh_paper_tables.sh
%  and then scripts/make_poster_figures.py.
%
%  Institutional logos: drop figures/logo-ceitec.pdf and figures/logo-muni.pdf
%  (PDF/EPS, or PNG at >=300 dpi) beside the figures and they are picked up
%  automatically. Until then a typographic lockup stands in.
% =====================================================================
\documentclass[final]{article}

\usepackage[paperwidth=841mm,paperheight=1189mm,
            left=30mm,right=30mm,top=24mm,bottom=20mm]{geometry}
\usepackage{amsmath}
\usepackage{fontspec}
\usepackage{unicode-math}
\usepackage{xcolor}
\usepackage{graphicx}
\usepackage{booktabs}
\usepackage{array}
\usepackage{tcolorbox}
\usepackage{qrcode}
\usepackage{microtype}
\tcbuselibrary{skins}

\pagestyle{empty}
\setlength{\parindent}{0pt}
\setlength{\fboxsep}{0pt}

% ------------------------------------------------------------------ palette
% Data colours are exactly the manuscript's validated categorical set, so the
% poster's accents and its figures are the same three hues.
\definecolor{btblue}{HTML}{2A78D6}
\definecolor{btorange}{HTML}{EB6834}
\definecolor{btviolet}{HTML}{4A3AA7}
\definecolor{btink}{HTML}{14181F}
\definecolor{btmuted}{HTML}{5A5F66}
\definecolor{btnavy}{HTML}{102A43}
\definecolor{btrule}{HTML}{E2E4E7}
\definecolor{btpage}{HTML}{F4F5F7}
\definecolor{btwash}{HTML}{EDF2FA}

\pagecolor{btpage}

% ------------------------------------------------------------------- fonts
% Poppins (display) + Carlito (text). Both are free:
%   Poppins  -- https://fonts.google.com/specimen/Poppins
%   Carlito  -- fonts-crosextra-carlito / https://fonts.google.com/specimen/Carlito
% If either is missing the poster still builds, in TeX Gyre Heros throughout.
\IfFontExistsTF{Carlito}
  {\setmainfont{Carlito}[Ligatures=TeX]}
  {\setmainfont{TeX Gyre Heros}[Ligatures=TeX]}
\setmathfont{Latin Modern Math}
\IfFontExistsTF{Poppins-Medium.ttf}
  {\newfontfamily\dsp{Poppins}[Ligatures=TeX,
     UprightFont=*-Medium, BoldFont=*-Bold, ItalicFont=*-MediumItalic]
   \newfontfamily\dsplight{Poppins}[Ligatures=TeX, UprightFont=*-Light]}
  {\newfontfamily\dsp{TeX Gyre Heros}[Ligatures=TeX]
   \newfontfamily\dsplight{TeX Gyre Heros}[Ligatures=TeX]}
\newfontfamily\mono{DejaVu Sans Mono}[Scale=0.80]

\newcommand{\sz}[2]{\fontsize{#1}{#2}\selectfont}
\newcommand{\body}{\sz{21}{29}}
\newcommand{\bodysm}{\sz{18.5}{25}}
\newcommand{\bodyxs}{\sz{16}{21.5}}

\newcommand{\x}{\,\text{\dsp\textbf{×}}\,}
\newcommand{\arr}{\,\textrightarrow\,}

% --------------------------------------------------------------- components
% A card: white, hairline, soft shadow. The shadow is what makes the grid read
% as layered rather than as a wall of boxes.
\tcbset{
  card/.style={
    enhanced, arc=4mm, colback=white, colframe=btrule, boxrule=0.7pt,
    left=11mm, right=11mm, top=8mm, bottom=8mm, boxsep=0mm,
    before skip=0pt, after skip=0pt,
    drop fuzzy shadow={black!15},
    before upper={\raggedright},
  },
  note/.style={
    enhanced, arc=3mm, boxrule=0pt, leftrule=2.8mm,
    left=7mm, right=7mm, top=6.5mm, bottom=6.5mm, boxsep=0mm,
    before skip=0pt, after skip=0pt,
    before upper={\raggedright},
  },
}

% Accent rule + title. This is the only heading level on the poster.
\newcommand{\cardhead}[2][btblue]{%
  {\color{#1}\rule{22mm}{2mm}}\par\vspace{4.5mm}%
  {\dsp\bfseries\sz{35}{43}\color{btink}\raggedright\hyphenpenalty=10000
   \exhyphenpenalty=10000 #2\par}\vspace{6.5mm}%
}

\newcommand{\lead}[1]{{\dsp\bfseries\color{btink}#1}}
\newcommand{\src}[1]{{\bodyxs\color{btmuted}\raggedright #1\par}}

% Tables are \input from paper/generated/, so a regenerated table can never
% disagree with the manuscript.
\newcommand{\fittable}[1]{%
  \par\nointerlineskip
  \begingroup
  \setlength{\tabcolsep}{3.5mm}\renewcommand{\arraystretch}{1.34}
  \resizebox{\linewidth}{!}{\input{#1}}%
  \endgroup\par}

% Always close the paragraph: a bare \includegraphics leaves horizontal mode
% open and the next line of prose is set on the figure's own baseline.
\newcommand{\posfig}[1]{\par\nointerlineskip
  \includegraphics[width=\linewidth]{#1}\par}

% One big finding in the dark band, and its caption underneath.
\newcommand{\keystat}[3]{%
  \begin{minipage}[t]{#1}\vspace{0pt}\raggedright
    {\dsp\bfseries\sz{90}{90}\color{#2}#3\par}\vspace{5mm}
  \end{minipage}}
\newcommand{\statcap}[2]{%
  \begin{minipage}[t]{#1}\vspace{0pt}\raggedright
    {\bodysm\color{white}#2\par}\end{minipage}}

% Column stack. Fixed height with \vfill between cards so all three columns
% end flush at the footer regardless of how much prose each one holds;
% \colh is the one number to retune if content is added or removed.
\newlength{\colh}\setlength{\colh}{729mm}
\newenvironment{poscol}[1]
  {\begin{minipage}[t][\colh][s]{#1}\vspace{0pt}\raggedright}
  {\vspace{0pt}\end{minipage}}
\newcommand{\cardgap}{\vfill}

% =====================================================================
\begin{document}

% ------------------------------------------------------------------ header
% The MUNI|CEITEC lockup is a single wide artwork, so it gets its own line
% above the title rather than a column beside it.
\IfFileExists{figures/logo-muni-ceitec.pdf}
  {\hfill\includegraphics[height=25mm]{figures/logo-muni-ceitec.pdf}\par
   \vspace{9mm}}
  {\hfill{\dsp\bfseries\sz{40}{40}\color{btink}MUNI \textbar{} CEITEC}\par
   \vspace{9mm}}

{\dsp\bfseries\sz{62}{72}\color{btink}\raggedright
 BioTerm-Bench: A Harbor-Compatible Benchmark for Tracking\\
 Open-Model Progress on Terminal Bioinformatics Tasks\par}

\vspace{6mm}
{\dsp\sz{40}{50}\color{btblue}\raggedright
 Can an LLM agent do your routine command-line bioinformatics?\par}

\vspace{7mm}
{\color{btink}\rule{\linewidth}{1.2mm}}\\[5mm]
{\dsp\sz{29}{36}\color{btink}Petr Simecek \quad\textbullet\quad
 \textbf{Vojtech Bystry}}\enspace
{\dsp\bfseries\sz{19}{19}\color{btblue}PRESENTING}%
\hfill
{\dsplight\sz{24}{31}\color{btmuted}Central European Institute of Technology
 (CEITEC), Masaryk University, Brno, Czechia}

\vspace{10mm}

% ------------------------------------------------------- key findings band
\begin{tcolorbox}[enhanced, arc=4mm, colback=btnavy, colframe=btnavy,
                  left=14mm, right=14mm, top=9mm, bottom=9mm, boxsep=0mm,
                  before skip=0pt, after skip=0pt, drop fuzzy shadow={black!22}]
  \keystat{0.225\linewidth}{white}{32}\hfill
  \keystat{0.225\linewidth}{white}{14}\hfill
  \keystat{0.225\linewidth}{btblue}{92\%}\hfill
  \keystat{0.21\linewidth}{btorange}{44\%}

  \statcap{0.225\linewidth}{self-contained terminal tasks, each graded by
    \textbf{comparison with a stored correct answer}}\hfill
  \statcap{0.225\linewidth}{models evaluated on both parts of the suite,
    across \textbf{two independent infrastructures}}\hfill
  \statcap{0.225\linewidth}{of individual attempts solved by \textbf{Gemma 4
    31B} --- a model that fits a laptop at four bits}\hfill
  \statcap{0.21\linewidth}{of failures are well-formed files carrying
    \textbf{wrong numbers} --- the mode nothing downstream catches}
\end{tcolorbox}

\vspace{11mm}

% ================================================================== columns
\newlength{\colw}\setlength{\colw}{0.3153\linewidth}

% ------------------------------------------------------------------ COL 1
\begin{poscol}{\colw}

\begin{tcolorbox}[card]
  \cardhead{The question}
  {\body
  Much of bioinformatics is not discovery. It is converting a GTF to a BED,
  checking that a BAM is coordinate-sorted, normalising counts to TPM ---
  precisely specified, tedious, easy to get subtly wrong.\par
  \vspace{5mm}
  \lead{Which parts can be delegated to an agent, and how often is the
  result right?}\par
  \vspace{5mm}
  Answering that needs a benchmark that runs the tools rather than describing
  them, grades deterministically rather than with an LLM judge, and runs on a
  laptop or an institutional node --- because the data that most needs help
  often cannot leave the building. Human genomic data under controlled access
  generally cannot be sent to public generative-AI services, so the agent has
  to run where the data already sits.\par}
\end{tcolorbox}

\cardgap

\begin{tcolorbox}[card]
  \cardhead[btviolet]{What BioTerm-Bench is}
  {\body
  \textbf{32 deterministic tasks} --- a 25-task core plus a 7-task extension of
  multi-stage workflows. Each hands the agent real input files, a container
  with the usual stack, and an \textbf{exact output contract}.\par
  \vspace{5mm}
  Grading compares the produced artifact with a stored correct answer.
  \lead{No model ever judges another model's output.}\par}
  \vspace{8mm}

  {\bodysm
  \setlength{\tabcolsep}{0pt}\renewcommand{\arraystretch}{1.42}
  \begin{tabular}{@{}p{0.34\linewidth}p{0.66\linewidth}@{}}
    \toprule
    alignment       & BAM depth, duplicate marking     \\
    variant         & allele frequencies, trio Mendelian \\
    metadata        & GTF\arr BED, longest ORF          \\
    quality control & per-base quality, FastQC summaries \\
    search          & best BLAST hit, \emph{k}-mer Jaccard \\
    RNA-seq         & TPM normalisation, DE contrasts   \\
    \bottomrule
  \end{tabular}\par}
  \vspace{8mm}

  {\body Runs under \textbf{Harbor}, the Terminal-Bench harness, with the agent
  held fixed at \textbf{Terminus-2}, so score differences reflect the model and
  not the scaffolding. Four attempts per task, and every fixture is small
  enough to rebuild and rerun on a laptop.\par}
\end{tcolorbox}

\cardgap

\begin{tcolorbox}[card]
  \cardhead[btorange]{The 25-task core is solved --- and that is recent}
  {\body \textbf{Gemma 4 31B} --- hostable on a laptop at four bits --- gets
  \textbf{92\,\%} of individual attempts right, and three models in the wider
  panel answer all 25 correctly on every scored attempt.\par}
  \vspace{8mm}
  \fittable{generated/poster_table_headline.tex}
  \vspace{6mm}
  \src{OpenRouter cohort, four attempts per task, 100 trials per model.
    ``all scored'' = solved on \emph{every} scored attempt; all denominators
    are explicit scored-task counts.}
  \vspace{7mm}
  {\body \lead{This says more about the moment than about the tasks.} Every
  model in the panel was released within the last twelve months; the same 25
  tasks run eighteen months ago would not have produced this table, and nine of
  the 32 are now failed by nobody at all.\par
  \vspace{5mm}
  The core keeps its value as a \textbf{continuity score} --- a record of what
  became routinely automatable, and how fast --- while the harder pack carries
  the discriminative load.\par}
\end{tcolorbox}

\end{poscol}%
\hfill
% ------------------------------------------------------------------ COL 2
\begin{poscol}{\colw}

\begin{tcolorbox}[card]
  \cardhead{Difficulty lives in reliability,\\not in solvability}
  \posfig{figures/pfig_paired_scopes.pdf}
  \vspace{4mm}
  \src{Every model evaluated on both parts of the suite; one row per model.}
  \vspace{8mm}
  \begin{tcolorbox}[note, colback=btwash, colframe=btblue]
    {\body\lead{Best-of-four hides the gap.} All five models of the OpenRouter
    cohort solve \emph{every} extension task within four attempts, so
    best-of-four calls them identical. Under the \textbf{all-four} criterion ---
    the one that matters when nobody re-runs the job --- the same five span
    \textbf{2 of 7 to 7 of 7} --- and a laboratory running each once gets a
    correct artifact on 100\,\% of attempts from one and 57\,\% from another.\par}
  \end{tcolorbox}
\end{tcolorbox}

\cardgap

\begin{tcolorbox}[card]
  \cardhead[btorange]{How agents fail --- and why it matters}
  \posfig{figures/pfig_failure_modes.pdf}
  \vspace{4mm}
  \src{All 319 scored failures in the included core cohort, classified from
    the verifiers' own reasons.}
  \vspace{8mm}
  \begin{tcolorbox}[note, colback=btorange!7, colframe=btorange]
    {\body\lead{The largest category is the quiet one.} A BED file with correct
    syntax and wrong coordinates, or a flagstat summary with plausible but
    incorrect counts, propagates silently into whatever consumes it. The loud
    56\% announce themselves --- a downstream tool refuses to read them. The
    quiet 44\% do not, and they are spread across ordinary tasks rather than
    exotic ones.\par}
  \end{tcolorbox}
\end{tcolorbox}

\end{poscol}%
\hfill
% ------------------------------------------------------------------ COL 3
\begin{poscol}{\colw}

\begin{tcolorbox}[card]
  \cardhead[btviolet]{What hardware do you actually need?}
  \posfig{figures/pfig_hardware_vs_score.pdf}
  \vspace{5mm}
  \fittable{generated/poster_table_hardware.tex}
  \vspace{6mm}
  \src{Weight memory is total parameters at the served precision --- a floor
    that excludes KV cache and runtime overhead. Letters mark the tier; tier A
    assumes 4-bit quantisation.}
  \vspace{6mm}
  {\body A \textbf{62\,GB workstation} model outscores a 122\,GB one needing a
  datacentre GPU.\par}
\end{tcolorbox}

\cardgap

\begin{tcolorbox}[card]
  \cardhead{Price is not capability}
  \posfig{figures/pfig_cost_capability.pdf}
  \vspace{6mm}
  {\body The whole extension sweep cost \textbf{\$0.13--\$1.64} per model, and
  the cheapest run ties for the best score. Price follows parameter count and
  provider margin --- not the licence.\par}
\end{tcolorbox}

\cardgap

\begin{tcolorbox}[card]
  \cardhead[btviolet]{Research infrastructure is the enabling layer}
  {\body \textbf{Nine of the fourteen models here were served by e-INFRA CZ},
  which offers LLM inference to \textbf{ELIXIR CZ} members on institutional
  hardware. That is what makes agentic analysis possible for controlled-access
  human data: the model runs \emph{inside the administrative boundary} and
  nothing is exported.\par
  \vspace{5mm}
  It is also sufficient. A tier B or C model hosted this way reaches the
  \textbf{top of the observed score range} --- no frontier API required.\par}
\end{tcolorbox}

\end{poscol}

\vspace{11mm}

% ------------------------------------------------------------------ footer
\begin{tcolorbox}[card, left=14mm, right=14mm, top=11mm, bottom=11mm]
  {\dsp\bfseries\sz{35}{43}\color{btink}What to delegate today\par}
  \vspace{8mm}

  \newcommand{\adv}[3]{%
    \begin{minipage}[t]{0.288\linewidth}\vspace{0pt}\raggedright
      {\dsp\bfseries\sz{25}{32}\color{#1}#2\par}\vspace{4mm}
      {\body\color{btink}#3\par}
    \end{minipage}}

  \adv{btblue}{Delegate, with spot-checks}{Format conversion and summarisation
    --- GTF\arr BED, VCF\arr TSV, FASTA statistics, flagstat-style counting.
    Reliable from the strongest models --- but also where weaker ones fail
    most, so check the artifact.}\hfill
  \adv{btviolet}{Delegate, then verify}{Tasks solved most of the time and
    plausibly wrong otherwise. Run it twice and diff --- cheap insurance
    against the failure mode nothing downstream will flag.}\hfill
  \adv{btorange}{Not yet}{Multi-stage reasoning over biological structure ---
    splice-junction support, trio-aware and diploid-consensus variant logic.
    One such task was solved by only 2 of 14 models.}

  \vspace{9mm}
  {\color{btrule}\rule{\linewidth}{0.8mm}}\\[8mm]

  \begin{minipage}[t]{0.695\linewidth}\vspace{0pt}\raggedright
    {\body\lead{Rule of thumb:} delegation is safe in proportion to how
    mechanically checkable the output is --- a BED file can be checked by
    re-deriving a handful of rows, a variant set with a subtle inheritance
    error cannot be checked without redoing the analysis.\par
    \vspace{5mm}
    \lead{And report mean reward and all-attempts, not best-of-\emph{n}.}
    Best-of-\emph{n} describes a workflow with an oracle telling you which run
    was right. You do not have one --- that is why you wanted the agent.\par}
    \vspace{7mm}
    {\bodysm\color{btmuted}Computational resources, and the served models
      themselves, were provided by the \textbf{e-INFRA CZ} project (ID:90254),
      supported by the Ministry of Education, Youth and Sports of the Czech
      Republic. Tasks, images, graders, runners and every analysis script are
      released under Apache~2.0.\par}
  \end{minipage}%
  \hfill
  \begin{minipage}[t]{0.265\linewidth}\vspace{0pt}
    \raggedleft
    {\dsp\bfseries\sz{18}{24}\color{btink}%
     github.com/ML-Bioinfo-CEITEC/BioTerm-Bench\par}
    \vspace{6mm}
    \qrcode[height=50mm,level=M]{https://github.com/ML-Bioinfo-CEITEC/BioTerm-Bench}
  \end{minipage}
\end{tcolorbox}

\end{document}
